\documentclass{amsart}
\usepackage{amsmath,amssymb}

\begin{document}

%% --- Section 1: math.align-columns targets ---

% cases with misaligned & columns
\begin{equation}
f(x) =
\begin{cases}
x^2 & \text{if } x > 0, \\
-x^3 + 1 & \text{if } x \le 0, \\
0 & \text{otherwise}.
\end{cases}
\end{equation}

% pmatrix with misaligned columns
The matrix is
\[
A = \begin{pmatrix} 1 & 0 & 0 \\ 0 & \alpha & \beta \\ 0 & \gamma & \delta \end{pmatrix}.
\]

% tabular with misaligned columns
\begin{tabular}{lll}
Short & Medium text & Long column \\
Longer entry & X & Y \\
A & B & C
\end{tabular}

% align* with misaligned columns
\begin{align*}
f(x) &= ax^2 + bx + c \\
g(x) &= \alpha x^3 + \beta x^2 \\
h(x) &= 1
\end{align*}

%% --- Section 2: math.continuation-indent targets ---

% equation with continuation lines starting with binops (not indented)
\begin{equation*}
S(n) = \sum_{k=1}^{n} k^2
+ \sum_{k=1}^{n} k
- n(n+1)/2.
\end{equation*}

% gather with continuation
\begin{gather*}
F(x) = \int_0^x f(t)\, dt
+ C_1
- C_2.
\end{gather*}

%% --- Section 3: prose.tilde-refs targets ---

We refer to Theorem 1.1 in \cite{ref1} and also \ref{thm:main}.
See equation \eqref{eq:first} for details.
Compare with \cref{sec:intro}.

%% --- Section 4: no-op environments (already aligned) ---

\begin{align*}
a &= b \\
c &= d
\end{align*}

%% --- Section 5: refusal cases for align-columns ---

% Unequal cell counts: should not align
% mreview-fmt: skip
\begin{align*}
a &= b &= c \\
d &= e
\end{align*}

% Comment in row: should not align
% mreview-fmt: skip
\begin{align*}
a &= b \\ % this is a comment
c &= d
\end{align*}

\end{document}
